% GENERATED FILE. Edit canonical lessons, then run npm run generate:books.
% canonical-source-hash: fnv1a64:ab09bd219c8adeb8
% canonical-lessons: GE-C463-raureif, GE-C463-gefrieren, GE-C463-tauen, GE-C463-tauwetter, GE-C463-bibbern

\chapter{Eisige Kälte}
\label{ch:ge-eisige-kalte}
\hlchaptermodality{463}

\begin{chapteropening}
\textbf{By the end of this chapter:} \emph{I can talk about frost and thaw: hoarfrost, freezing, thawing, the thaw and shivering.}

Complete the last lesson of chapter 463: I can talk about frost and thaw: hoarfrost, freezing, thawing, the thaw and shivering.
\end{chapteropening}

\section[der Raureif]{der Raureif — the hoarfrost}

\label{lesson:GE-C463-raureif}



\begin{quote}
Before the new one: say the German for to hail, then the German for to thunder.
\end{quote}

\subsection*{You'll want to know: der Raureif}
\textbf{der Raureif} — \textquotedblleft{}the hoarfrost\textquotedblright{}.

\begin{quote}
Am Morgen sind die Bäume weiß vom Raureif.
\end{quote}

\begin{quote}
\emph{In the morning the trees are white with hoarfrost.}
\end{quote}

No plural. Also simply \emph{der Reif}.

\begin{cousinweb}
\emph{Raureif} is built from \emph{rau} 'rough' and \emph{Reif} 'hoarfrost', which EtymDB traces through Middle High German \emph{rīfe} to Old High German \emph{rīfo}. This \emph{Reif} is a different word from \emph{Reif} 'ring', though EtymDB lists both under one entry.
\end{cousinweb}

\begin{grammarlens}[title={What you've built}]
One more word for telling, explaining and describing at B1.
\end{grammarlens}

\subsection*{Your turn}
\begin{itemize}
\raggedright
  \item \emph{Say it:} \emph{der Raureif}
  \item \emph{Say it:} \emph{der Raureif}, once more
  \item \emph{Say it:} a sentence of your own with \emph{der Raureif}
  \item \emph{Recall:} say the German for to hail, then the German for to thunder, then say \emph{der Raureif} again
\end{itemize}

\begin{tcolorbox}[breakable,colback=teal!4,colframe=teal!35!black,title={Before you move on}]
What does der Raureif mean? (\textquotedblleft{}The hoarfrost\textquotedblright{}.) Say it once more.
\end{tcolorbox}
\section[gefrieren]{gefrieren — to freeze (solid)}

\label{lesson:GE-C463-gefrieren}



\begin{quote}
Before the new one: say the German for to thunder, then the German for the hoarfrost.
\end{quote}

\subsection*{You'll want to know: gefrieren}
\textbf{gefrieren} — \textquotedblleft{}to freeze (solid)\textquotedblright{}.

\begin{quote}
Bei minus zehn Grad gefriert sogar der See.
\end{quote}

\begin{quote}
\emph{At minus ten degrees even the lake freezes.}
\end{quote}

Perfekt with \emph{sein}: \emph{der See ist gefroren}.

\begin{cousinweb}
\emph{gefrieren} is built from \emph{ge-} and the inherited verb for freezing; EtymDB traces \emph{frieren} through Old High German \emph{friosan} to Proto-Germanic \emph{freusaną} and lists Dutch \emph{vriezen} and English \emph{freeze} as cousins.
\end{cousinweb}

\begin{grammarlens}[title={What you've built}]
One more word for telling, explaining and describing at B1.
\end{grammarlens}

\subsection*{Your turn}
\begin{itemize}
\raggedright
  \item \emph{Say it:} \emph{gefrieren}
  \item \emph{Say it:} \emph{gefrieren}, once more
  \item \emph{Say it:} a sentence of your own with \emph{gefrieren}
  \item \emph{Recall:} say the German for to thunder, then the German for the hoarfrost, then say \emph{gefrieren} again
\end{itemize}

\begin{tcolorbox}[breakable,colback=teal!4,colframe=teal!35!black,title={Before you move on}]
What does gefrieren mean? (\textquotedblleft{}To freeze (solid)\textquotedblright{}.) Say it once more.
\end{tcolorbox}
\section[tauen]{tauen — to thaw, to melt}

\label{lesson:GE-C463-tauen}



\begin{quote}
Before the new one: say the German for the hoarfrost, then the German for to freeze.
\end{quote}

\subsection*{You'll want to know: tauen}
\textbf{tauen} — \textquotedblleft{}to thaw, to melt\textquotedblright{}.

\begin{quote}
Im März taut der Schnee auf den Dächern.
\end{quote}

\begin{quote}
\emph{In March the snow on the roofs thaws.}
\end{quote}

With \emph{es} for the weather: \emph{es taut}. Perfekt with \emph{sein} for snow and ice: \emph{der Schnee ist getaut}.

\begin{cousinweb}
\emph{tauen} 'to thaw' is an inherited Germanic verb (EtymDB lists 'to thaw'), and English \emph{thaw} is a cousin. It is a different word from \emph{Tau} 'dew'.
\end{cousinweb}

\begin{grammarlens}[title={What you've built}]
One more word for telling, explaining and describing at B1.
\end{grammarlens}

\subsection*{Your turn}
\begin{itemize}
\raggedright
  \item \emph{Say it:} \emph{tauen}
  \item \emph{Say it:} \emph{tauen}, once more
  \item \emph{Say it:} a sentence of your own with \emph{tauen}
  \item \emph{Recall:} say the German for the hoarfrost, then the German for to freeze, then say \emph{tauen} again
\end{itemize}

\begin{tcolorbox}[breakable,colback=teal!4,colframe=teal!35!black,title={Before you move on}]
What does tauen mean? (\textquotedblleft{}To thaw, to melt\textquotedblright{}.) Say it once more.
\end{tcolorbox}
\section[das Tauwetter]{das Tauwetter — the thaw}

\label{lesson:GE-C463-tauwetter}



\begin{quote}
Before the new one: say the German for to freeze, then the German for to thaw.
\end{quote}

\subsection*{You'll want to know: das Tauwetter}
\textbf{das Tauwetter} — \textquotedblleft{}the thaw\textquotedblright{}.

\begin{quote}
Mit dem Tauwetter werden die Wege nass und matschig.
\end{quote}

\begin{quote}
\emph{With the thaw the paths become wet and muddy.}
\end{quote}

No plural. Also figurative, in politics.

\begin{cousinweb}
\emph{Tauwetter} is built from \emph{tauen} 'to thaw' and \emph{Wetter} 'weather', both inherited Germanic words.
\end{cousinweb}

\begin{grammarlens}[title={What you've built}]
One more word for telling, explaining and describing at B1.
\end{grammarlens}

\subsection*{Your turn}
\begin{itemize}
\raggedright
  \item \emph{Say it:} \emph{das Tauwetter}
  \item \emph{Say it:} \emph{das Tauwetter}, once more
  \item \emph{Say it:} a sentence of your own with \emph{das Tauwetter}
  \item \emph{Recall:} say the German for to freeze, then the German for to thaw, then say \emph{das Tauwetter} again
\end{itemize}

\begin{tcolorbox}[breakable,colback=teal!4,colframe=teal!35!black,title={Before you move on}]
What does das Tauwetter mean? (\textquotedblleft{}The thaw\textquotedblright{}.) Say it once more.
\end{tcolorbox}
\section[bibbern]{bibbern — to shiver, to tremble (colloquial)}

\label{lesson:GE-C463-bibbern}



\begin{quote}
Before the new one: say the German for to thaw, then the German for the thaw.
\end{quote}

\subsection*{You'll want to know: bibbern}
\textbf{bibbern} — \textquotedblleft{}to shiver, to tremble (colloquial)\textquotedblright{}.

\begin{quote}
An der Haltestelle bibbern die Kinder vor Kälte.
\end{quote}

\begin{quote}
\emph{At the bus stop the children shiver with cold.}
\end{quote}

Everyday word, with \emph{vor}: \emph{vor Kälte bibbern}.

\begin{cousinweb}
\emph{bibbern} is an expressive form usually linked to \emph{beben} 'to tremble'; EtymDB does not list it, so the line is hedged.
\end{cousinweb}

\begin{grammarlens}[title={What you've built}]
One more word for telling, explaining and describing at B1.
\end{grammarlens}

\subsection*{Your turn}
\begin{itemize}
\raggedright
  \item \emph{Say it:} \emph{bibbern}
  \item \emph{Say it:} \emph{bibbern}, once more
  \item \emph{Say it:} a sentence of your own with \emph{bibbern}
  \item \emph{Recall:} say the German for to thaw, then the German for the thaw, then say \emph{bibbern} again
\end{itemize}

\begin{tcolorbox}[breakable,colback=teal!4,colframe=teal!35!black,title={Before you move on}]
What does bibbern mean? (\textquotedblleft{}To shiver, to tremble (colloquial)\textquotedblright{}.) Say it once more.
\end{tcolorbox}
